\begin{enumerate}

\item El centro del círculo es $30$ MPa.

\item El radio es $\sqrt{50^2 + 30^2}$.

\item Las tensiones principales son el centro más y menos el radio; el radio es la tensión tangencial máxima.

\item El ángulo se obtiene de $\tan 2\theta_p = 2\tau_{xy}/(\sigma_x - \sigma_y)$.

\end{enumerate}

\noindent\textbf{Operaciones y resultados intermedios.}

\begin{enumerate}

\item Radio del círculo de Mohr.

\[\fitmath{y_{1}=\sqrt{{\left(\frac{80--20}{2}\right)}^{2}+{30}^{2}}\approx 58{,}30952\,\mathrm{MPa}}\]

\item Tensión principal menor.

\[\fitmath{y_{2}=\frac{80+-20}{2}-\sqrt{{\left(\frac{80--20}{2}\right)}^{2}+{30}^{2}}\approx -28{,}30952\,\mathrm{MPa}}\]

\item Ángulo del plano principal en grados.

\[\fitmath{y_{3}=\frac{\frac{\arctan\left(\frac{2\cdot 30}{80--20}\right)}{2}\cdot 180}{\pi}\approx 15{,}48188\,{}^{\circ}}\]

\item Tensión principal mayor.

\[\fitmath{y_{4}=\frac{80+-20}{2}+\sqrt{{\left(\frac{80--20}{2}\right)}^{2}+{30}^{2}}\approx 88{,}30952\,\mathrm{MPa}}\]

\end{enumerate}
